\section*{Hoofdstuk \ref{ch:g12:catalysis} --- Katalyse}

\begin{solution}{exo:g12:catalysis:1}
\omterm{def:g7:chemical-reactions:reactant}{Beginstof}: \ce{H2O2}. \omterm{def:g7:chemical-reactions:reactant}{Reactieproducten}: \ce{H2O} en \ce{O2}. \omterm{def:g12:catalysis:catalyst}{Katalysator}:
\ce{I-} (gebruikt in stap 1, teruggegeven in stap 2). Intermediair:
\ce{IO-} (gevormd in stap 1, gebruikt in stap 2).
\end{solution}

\begin{solution}{exo:g12:catalysis:2}
Homogeen; heterogeen; heterogeen; homogeen.
\end{solution}

\begin{solution}{exo:g12:catalysis:3}
De niet-gekatalyseerde weg. Het begin- en eindniveau zijn dezelfde:
dezelfde \omterm{def:g7:chemical-reactions:reactant}{reactieproducten} ontstaan, met dezelfde vrijgekomen energie;
alleen de weg verschilt.
\end{solution}

\begin{solution}{exo:g12:catalysis:4}
Hij wordt in de ene stap gebruikt en in een andere teruggegeven: in de som
van de stappen staat hij aan beide kanten, en valt hij weg.
\end{solution}

\begin{solution}{exo:g12:catalysis:5}
De reactie vindt aan het oppervlak plaats: hoe groter het oppervlak, hoe
meer \omterm{def:g7:atoms-and-molecules:molecule}{moleculen} tegelijk kunnen reageren, voor dezelfde massa (vaak dure)
\omterm{def:g12:catalysis:catalyst}{katalysator}.
\end{solution}

\begin{solution}{exo:g12:catalysis:6}
Som: \ce{2Fe^{3+} + 2H2O2 + 2Fe^{2+} + 2H+ -> 2Fe^{2+} + O2 + 2H+ +
2Fe^{3+} + 2H2O}. De ijzerionen en de waterstofionen staan aan beide
kanten en vallen weg: \ce{2H2O2 -> 2H2O + O2}.
\end{solution}

\begin{solution}{exo:g12:catalysis:7}
De helft van \qty{72}{mL} is \qty{36}{mL}: bereikt na ongeveer
\qty{35}{min} zonder \omterm{def:g12:catalysis:catalyst}{katalysator} en \qty{3.5}{min} met \omterm{def:g12:catalysis:catalyst}{katalysator}. De
\omterm{def:g12:catalysis:catalyst}{katalysator} maakt die tijd tien keer korter.
\end{solution}

\begin{solution}{exo:g12:catalysis:8}
De reactie vindt plaats aan het oppervlak van het platina; het lood
bedekt dat en blijft zitten, zodat de gassen het \omterm{def:g9:periodic-table-first-look:metal}{metaal} niet meer
bereiken. De \omterm{def:g12:catalysis:catalyst}{katalysator} is ``vergiftigd'': aanwezig maar nutteloos.
\end{solution}

\begin{solution}{exo:g12:catalysis:9}
Tot ongeveer \qty{37}{\celsius} versnelt de reactie met de temperatuur,
zoals elke reactie. Daarboven verliest het \omterm{def:g12:catalysis:enzyme}{enzym}, een eiwit, zijn vorm en
daarmee zijn activiteit; bij \qty{80}{\celsius} is het vernietigd.
\end{solution}

\begin{solution}{exo:g12:catalysis:10}
\ce{2CO + 2NO -> 2CO2 + N2}: C 2 = 2, O 4 = 4, N 2 = 2. Koolstofmono-oxide
wordt geoxideerd door stikstofmono-oxide, dat gereduceerd wordt: elke
verontreiniging ruimt de andere op.
\end{solution}

\begin{solution}{exo:g12:catalysis:11}
$n(\ce{O2}) = 0.072 / 24.1 = \qty{2.99e-3}{mol}$;
$n(\ce{H2O2}) = 2 \times \num{2.99e-3} = \qty{5.98e-3}{mol}$ in
\qty{10.0}{mL}: $c = \qty{0.598}{mol/L}$.
\end{solution}

\begin{solution}{exo:g12:catalysis:12}
\ce{C2H5OH -> C2H4 + H2O}, een \omterm{def:g12:curly-arrows:categories}{eliminatie}; \ce{C2H5OH -> CH3CHO + H2}.
Een \omterm{def:g12:catalysis:selective}{selectieve katalysator} versnelt een van meerdere mogelijke reacties
veel meer dan de andere, en bepaalt zo het product.
\end{solution}

\begin{solution}{exo:g12:catalysis:13}
De \omterm{def:g12:catalysis:catalyst}{katalysator} verandert de \omterm{def:g10:reaction-progress-table:system}{eindtoestand} niet: de krommen met en zonder
\omterm{def:g12:catalysis:catalyst}{katalysator} eindigen bij hetzelfde volume zuurstof. Je krijgt dezelfde
hoeveelheid product, alleen eerder.
\end{solution}

\begin{solution}{exo:g12:catalysis:14}
$\num{6.02e23} / \num{5e6} = \qty{1.2e17}{s}$, zo'n vier miljard jaar.
Om het in één seconde te doen: \num{1.2e17} enzymmoleculen, dat is
\qty{2e-7}{mol}.
\end{solution}

\begin{solution}{exo:g12:catalysis:15}
Ammoniak \ce{NH3}: $150 \times 17.0 / 14.0 = \qty{182}{Mt}$.
Stikstofatomen: $\num{1.50e14} / 14.0 = \qty{1.07e13}{mol}$, dat is
\qty{5.36e12}{mol} \ce{N2}.
\end{solution}

\begin{solution}{pb:g12:catalysis:1}
\textbf{1.} \ce{2H2O2 -> 2H2O + O2}.

\textbf{2.} Ja: waterstofperoxide is de \omterm{def:g11:redox:oxidant}{oxidator} van het koppel
\ce{H2O2}/\ce{H2O} en de \omterm{def:g11:redox:oxidant}{reductor} van het koppel \ce{O2}/\ce{H2O2}; het
oxideert en reduceert zichzelf.

\textbf{3.} Zonder \omterm{def:g12:catalysis:catalyst}{katalysator} is de reactie extreem langzaam.

\textbf{4.} $20 / 22.4 = \qty{0.893}{mol}$.

\textbf{5.} Twee \omterm{def:g10:the-mole:amount}{mol} waterstofperoxide per \omterm{def:g10:the-mole:amount}{mol} zuurstof:
\qty{1.79}{mol}.

\textbf{6.} \qty{1.79}{mol/L}.

\textbf{7.} $M = 2 \times 1.0 + 2 \times 16.0 = \qty{34.0}{g/mol}$;
$C_m = 1.79 \times 34.0 = \qty{60.7}{g/L}$.

\textbf{8.} De helft: \qty{0.893}{mol/L}.

\textbf{9.} Katalase: een \omterm{def:g12:catalysis:enzyme}{enzym}, opgelost, een biologische \omterm{def:g12:catalysis:catalyst}{katalysator};
ijzer(III)\omterm{def:g9:ions:ion}{ionen}: \omterm{def:g12:catalysis:homogeneous}{homogene katalyse}.

\textbf{10.} \ce{2Fe^{3+} + H2O2 -> 2Fe^{2+} + O2 + 2H+} en
\ce{2Fe^{2+} + H2O2 + 2H+ -> 2Fe^{3+} + 2H2O}; hun som is
\ce{2H2O2 -> 2H2O + O2}.

\textbf{11.} Met een \omterm{def:g12:catalysis:catalyst}{katalysator} stijgt het volume veel sneller; beide
krommen vlakken af bij hetzelfde eindvolume.

\textbf{12.} De ijzer(III)\omterm{def:g9:ions:ion}{ionen} (oranje) worden in de eerste stap
ijzer(II) (lichtgroen) en in de tweede weer ijzer(III); aan het eind is
al het ijzer weer ijzer(III).

\textbf{13.} Het koken heeft het \omterm{def:g12:catalysis:enzyme}{enzym} vernietigd.

\textbf{14.} $0.100 \times 20 = \qty{2.0}{L}$.

\textbf{15.} $n(\ce{O2}) = 0.100 \times 0.893 = \qty{0.0893}{mol}$, dat
is $0.0893 \times 24.1 = \qty{2.15}{L}$.

\textbf{16.} Door de langzame ontleding komt zuurstof vrij, die in een
goed afgesloten fles de druk zou opbouwen.

\textbf{17.} $1.50 / 2 = \qty{0.750}{mol}$ zuurstof per liter, dat is
$0.750 \times 22.4 = \qty{16.8}{L}$: ``16.8 volumes''.

\textbf{18.} $(1.79 - 1.50)/1.79 \approx \qty{16}{\percent}$.

\textbf{19.} \qty{1.79}{mol/L}.
\end{solution}
